\scanpage{14}\mspage{\handwritten{20}}
\[
  \handwritten{\mu}^{i}_{y},\quad{\handwritten{\nu}}^{i}_{y}
  \tag*{(4.1\deleted{7}\handwritten{9})}
\]
the radicial multiplicity and the separable multipli\makebox[0pt][l]{\raisebox{-.72em}[0pt][0pt]{\scriptsize;}}city of $C^{i}_{\deleted{i}\handwritten{y}}$ over $k(y)$ respectively
\reversemarginpar\margin{\handwritten{\unsure{x|$\times$}}}\normalmarginpar
(EGA IV 4\handwritten{.7.4.} ), - thus $\handwritten{\mu}^{i}_{y}=1$ if $k(y)$ is per\deleted{a}fect.

We \deleted{also set} also consider the smallest Galois extension $k(y)^{i}$
\figure{\begin{tikzpicture}[x=1mm,y=1mm,line cap=round,every node/.style={anchor=base west,inner sep=0pt,font=\scriptsize}]
  \node at (0,15) {$\deleted{(4.18)}$};
  \node at (35,15) {$\deleted{Z^{i}_{y}=\operatorname{Spec}(H^{0}(X^{i}_{y\,\mathrm{réd}},\underline O_{X^{i}_{y\,\mathrm{réd}}}))\,,\qquad Z=\unsure{\textstyle\prod|\textstyle\coprod}_{i}Z^{i}_{y}\,,}$};
  \node at (0,7) {\deleted{so that we simply have $Z^{i}_{y}$ is the spectrum of a finite extension}};
  \node at (0,0) {\deleted{$k(y)_{i}$ of $k(y)$, of separable degree}};
  \draw (0,-.5) -- (0,16) .. controls (10,17) and (26,18) .. (40,18)
    -- (146,17) .. controls (151,17) and (151,15) .. (155,11);
  \draw (154,10) -- (72,8) -- (72,-1) .. controls (50,-2) and (22,-1) .. (0,-1);
  \draw (18,15) .. controls (21,9) and (27,3) .. (33,0);
  \draw (42,15) .. controls (46,12) and (49,5) .. (53,3);
  \draw (66,15) .. controls (67,11) and (71,9) .. (77,8);
  \draw (91,15) .. controls (93,13) and (94,11) .. (97,10);
  \draw (109,12) -- (114,10);
  \draw (20,11.5) -- (20.4,11.1);
\end{tikzpicture}} of $k(y)$ such that the irreducible components of $(X^{i}_{y})\otimes_{k(y)} k(y)^{i}$ are geometrically irreducible: it is a finite extension\tikz[baseline=0pt,x=1em,y=1ex,line cap=round]{
  \path[use as bounding box] (0,-.7) rectangle (.1,4.0);
  \node[overlay,anchor=base west,inner sep=0pt] at (-1.0,2.5) {\inserted{\handwritten{of $k(y)$,}}};
  \draw[overlay] (0,-.7) .. controls (-.3,1.2) and (.2,1.4) .. (.7,1.4)
    -- (5.1,1.6) .. controls (5.5,1.6) and (5.6,1.9) .. (5.6,2.7);
} of degree $\handwritten{\nu}^{i}_{y}$. We set
\reversemarginpar\margin{\tikz[x=1em,y=1ex,line cap=round]{
  \draw[gray] (0,1.3) .. controls (.5,.8) and (1.0,.5) .. (1.8,.25);
  \draw (.45,.45) .. controls (.7,.15) and (.85,.25) .. (.85,.5)
    .. controls (1.0,.25) and (1.3,.5) .. (1.7,.25);
}}\normalmarginpar
\[
  Z^{i}_{y}=\operatorname{Spec}(k(y)^{i}),\qquad
  Z_{\deleted{\unsure{x|a|i}}y}=\handwritten{\textstyle\prod}Z^{i}_{y}\quad\deleted{\unsure{,|;}}\handwritten{,}
  \tag*{(4.\deleted{18}\handwritten{20})}
\]
This being so, the sheaf $\underline E_{y}$ on $y$ fits into an exact sequence of sheaves
\[
  0\handwritten{\to}\underline{\mathrm Z}_{y}\handwritten{\longrightarrow}
  p_{y\handwritten{*}}(\underline{\mathrm Z}_{Z_{y}})\handwritten{\to}
  \underline E_{y}\handwritten{\longrightarrow}0\quad ,
  \tag*{(4.\deleted{19}\handwritten{21})}
\]
where $p_{y}:Z_{y}\handwritten{\longrightarrow}y$ is the structure projection, and where the homomorphism $\underline{\mathrm Z}_{y}\handwritten{\longrightarrow}p_{y\handwritten{*}}(\underline{\mathrm Z}_{Z_{y}})$ is adjoint to the homomorphism
\[
  p_{y}^{\handwritten{*}}(\underline{\mathrm Z}_{y})=\underline{\mathrm Z}_{Z}
  \handwritten{\longrightarrow}\underline{\mathrm Z}_{Z_{\deleted{y}}}
  \tag*{(4.2\deleted{0}\handwritten{2})}
\]
which, on the component $Z^{i}_{y}$, reduces to multiplication by $a^{i}_{y}$. The exact cohomology sequence associated with (4.\deleted{19}\handwritten{21}) then gives us
\[
\begin{aligned}[t]
  H^{1}(y,\underline E_{y})&=\deleted{\unsure{\mathrm{Co}}}\operatorname{Ker}
  \bigl(H^{2}(y,\underline{\mathrm Z}_{y})\handwritten{\longrightarrow\textstyle\prod_{i}}
  H^{2}(Z^{i}_{y},\underline{\mathrm Z})\bigr)\\
  &=\deleted{\unsure{\mathrm{Co}}}\operatorname{Ker}
  \bigl(H^{1}(y,\underline{\mathrm Q}/\underline{\mathrm Z})\handwritten{\to}
  \handwritten{\tikz[baseline=-.6ex]{
  \node[inner sep=0pt] (p) {$\displaystyle\prod_i$};
  \node[overlay,inner sep=0pt] at ([yshift=.3ex]p.center) {$\deleted{\unsure{\to}}$};
}}H^{1}(Z^{i}_{y},\underline{\mathrm Q}/\underline{\mathrm Z})\bigr)\quad ,
\end{aligned}
\tag*{(4.2\deleted{1}\handwritten{3})}
\]
where the arrow $H^{1}(y,\underline{\mathrm Q}/\underline{\mathrm Z})\handwritten{\longrightarrow}H^{1}(Z^{i}_{y},\underline{\mathrm Q}/\underline{\mathrm Z})$ is the obvious tra\deleted{o}\handwritten{n}\deleted{i}sition homomorphism, multiplied by $a^{i}_{y}$. Since $H^{1}(y,\underline E_{y})$ is the intersection of the kernels of these homomorphisms, \tikz[baseline=(old.base)]{
  \node[inner sep=0pt,anchor=base west] (old) {\deleted{each of which}};
  \node[overlay,inner sep=0pt,anchor=base west] at ([yshift=1.0em]old.base west) {\inserted{and since such a kernel}};
  \path (old.north) ++(0,1.0em);
} is manifestly annihilated by $a^{i}_{y}\handwritten{\nu}^{i}_{y}$\deleted{\unsure{e}}, we see that $H^{1}(y,\underline E_{y})$ is annihilated by
\[
  d_{y}=\operatorname{pgcd}(a^{i}_{y}\handwritten{\nu}^{i}_{y})\deleted{\unsure{;}}.
  \tag*{(4.2\deleted{2}\handwritten{4})}
\]

When $k(y)$ is finite, or more generally pseudo-finite in the sense of Serre, i.e. its fundamental group $\operatorname{Gal}(\overline{k(y)}/k(y))$ isomorphic to $\overset{\handwritten{\scriptscriptstyle\wedge}}{\underline{\mathrm Z}}$, then we f\deleted{o}ind
% archive sheet number: 13
\clearpage
